% Pista ana-25j-q1 · Anàlisi
% Procedència: PAU Catalunya 2025 · Convocatòria de juny · Sèrie 1 · Exercici 1
\documentclass[11pt,a4paper]{article}
\usepackage{pista}
\pistainfo{Catalunya 2025}{Convocatòria de juny}{Sèrie 1}

\begin{document}

\section*{\textcolor{blauPista}{Exercici 1}}
\addcontentsline{toc}{section}{Exercici 1}

\noindent Considereu la funció
\[
f(x) = \dfrac{x^2 - 2x}{x - 1}.
\]

\subsection*{\textit{a)} \normalfont Determineu els talls de la corba $y = f(x)$ amb els eixos de coordenades, i les equacions de les seves possibles asímptotes verticals, horitzontals i obliqües. \hfill\textnormal{[1~punt]}}

\pista{Troba les arrels del polinomi numerador, per trobar els punts de tall amb l'eix $OX$; en canvi, el punt de tall amb l'eix $OY$ és més immediat, perquè només cal avaluar $f(0)$. Per a les asímptotes verticals, mira on s'anul·la el denominador i comprova que el límit de $f$ hi és infinit. Per a les altres, compara el grau del numerador i del denominador: quan el grau del polinomi numerador és 1 unitat superior al grau del polinomi denominador, llavors sempre hi haurà una asímptota obliqua, $y = mx + n$, amb $m = \lim_{x\to\infty} \frac{f(x)}{x}$ i $n = \lim_{x\to\infty} \bigl(f(x) - mx\bigr)$.}

\subsection*{\textit{b)} \normalfont Calculeu les equacions de les rectes tangents a la corba $y = f(x)$ en els punts $x = 0$ i $x = 2$. Aquestes dues rectes són paral·leles? Justifiqueu la resposta. \hfill\textnormal{[1~punt]}}

\pista{Recorda la fórmula $y - f(a) = f'(a)(x - a)$. Per derivar $f(x)$, aplica la regla del quocient. Dues rectes són paral·leles si i només si tenen el mateix pendent: compara $f'(0)$ i $f'(2)$.}

\subsection*{\textit{c)} \normalfont Hi ha algun punt on la recta tangent a $f(x)$ tingui pendent 1? En cas afirmatiu, trobeu-lo. \hfill\textnormal{[0,5~punts]}}

\pista{El pendent de la tangent en un punt és la derivada en aquell punt. Planteja l'equació $f'(x) = 1$ i analitza si té solució dins el domini.}

\vspace{0.6em}

\end{document}
